\section{Regular names, shrinking viability, and guarded operations}
\label{sec:namesguards}

\subsection{Regular rational names}

A convenient constructive real presentation is a rational sequence
\(q:\N\to\Q\) with a fixed regularity law, for example
\[
       |q_{p+1}-q_p|\le 2^{-p}.
\]
The law implies a primitive-recursive tail estimate
\[
       |q_m-q_p|\le 2^{1-p}
       \qquad(m\ge p).
\]
All statements are finite rational inequalities.  A completed real is never
needed as a set-theoretic object inside iPRA.

More generally, a name may come with a modulus \(M(p)\) such that
\[
 m,n\ge M(p)\quad\Longrightarrow\quad |q_m-q_n|\le 2^{-p}.
\]
The modulus must be explicit.  Mere convergence or metastability is not enough
for operations requiring a uniform tail.

\subsection{Names from relational towers}

Let \(v(p,a)\in\Q\) be the value of candidate \(a\) at level \(p\).  For an
increasing primitive-recursive horizon \(H(p)\), choose the least candidate at
level \(p\) that lies on a viable path to \(H(p)\); call its value \(q_p\).
Assume:
\begin{align}
 \diam\{v(p,a):V_{H}(p,a)\}&\le\delta_p
   &&\text{for every }H\ge H(p),\label{eq:diambound}\\
 E_p(a,b)&\Longrightarrow |v(p,a)-v(p+1,b)|\le\mu_p,
   \label{eq:motionbound}
\end{align}
and there is a primitive-recursive tail bound
\[
       \sum_{j=p}^{m-1}\mu_j\le T(p)
       \quad(m>p),
       \qquad T(p)\to0
\]
in an explicit numerical form.

\begin{theorem}[Shrinking viability gives a name]
\label{thm:shrinkingname}
For \(m>p\),
\[
      |q_m-q_p|\le \delta_p+T(p)+\delta_m.
\]
Consequently, if the right-hand side is at most \(2^{-r}\) whenever
\(p,m\ge M(r)\), then \((q_p)\) is a Cauchy name with modulus \(M\).
\end{theorem}

\begin{proof}
Take a viable path witnessing the choice at horizon \(H(m)\).  Its level-\(p\)
candidate, call it \(a_p\), remains viable after truncating that path at
horizon \(H(p)\), because \(H(m)\ge H(p)\).  The candidate defining \(q_p\) is
also viable to \(H(p)\).  Apply \eqref{eq:diambound} at the common horizon
\(H(p)\) to conclude that \(q_p\) and \(v(p,a_p)\) differ by at most
\(\delta_p\).  Follow the longer path from level \(p\) to level \(m\);
telescoping \eqref{eq:motionbound} bounds the change by \(T(p)\).  At level
\(m\), the candidate on this path and the independently selected candidate
defining \(q_m\) are both viable to the common horizon \(H(m)\), so their
values differ by at most \(\delta_m\).  Add the three rational inequalities.
Every path involved is finite, so the proof is a bounded certificate argument.
\end{proof}

This theorem avoids selecting one infinite branch.  Each \(q_p\) uses a finite
horizon; shrinking diameter makes the horizon dependence observationally
irrelevant.

\subsection{Three-way guards}

Exact binary branching on approximate data is generally discontinuous.  Given
a rational approximation \(q\) to a quantity \(x\), an error bound
\(|x-q|\le e\), and a positive margin \(\eta\), define
\[
 \operatorname{guard}(q,e,\eta)=
 \begin{cases}
  -,&q+e< -\eta,\\
  +,&q-e> \eta,\\
  ?,&\text{otherwise}.
 \end{cases}
\]
The output is a finite tag computed by rational comparison.

\begin{lemma}[Guard soundness]
If the guard returns \(+\), then \(x>\eta\).  If it returns \(-\), then
\(x<-\eta\).  If it returns \(?\), then the approximation has not certified a
side; no conclusion is forced.
\end{lemma}

\begin{proof}
For the positive case, \(x\ge q-e>\eta\).  The negative case is symmetric.
The unknown case is deliberately information-preserving.
\end{proof}

As precision increases, the error budget decreases.  An unknown guard may
later resolve, but a certified positive or negative guard remains sound.  When
the theorem assumes an apartness margin \(|x|>\eta\), sufficiently strong
precision guarantees resolution.

\subsection{Guarded reciprocal}

Let a regular name \(q_p\) represent \(x\), with
\(|x-q_p|\le e_p\).  Suppose an explicit rational \(\eta>0\) satisfies
\(|x|\ge\eta\).  Choose \(p\) with \(e_p\le\eta/2\).  Then
\(|q_p|\ge\eta/2\), so rational division by \(q_p\) is safe.

For two approximations \(q,r\) with
\(|x-q|,|x-r|\le e\le\eta/2\),
\[
 \left|\frac1q-\frac1r\right|
 =\frac{|q-r|}{|qr|}
 \le \frac{2e}{(\eta/2)^2}
 =\frac{8e}{\eta^2}.
\]
Thus, to obtain reciprocal precision \(2^{-k}\), it suffices to request an
input error at most \(2^{-k}\eta^2/8\).  This yields a primitive-recursive
backward demand once \(\eta\) is coded rationally.

\begin{theorem}[Guarded reciprocal]
There is a contracted transducer taking a regular rational name for \(x\), a
positive rational \(\eta\), and a certificate \(|x|\ge\eta\), and returning a
regular name for \(1/x\).  Its demand transformer is primitive recursive and
its correctness proof uses only rational inequalities.
\end{theorem}

The same pattern handles matrix inversion from a determinant or singular-value
margin, Gaussian pivots from a pivot margin, chart selection from a distance-to-
boundary margin, logarithms on positive intervals, and isolated algebraic
branches.

\subsection{Root and eigenvalue labels}

At a collision, an ordered list of roots or eigenvectors is not continuous.
The correct unconditional outputs are:
\begin{itemize}
\item a multiset or cluster cloud;
\item symmetric functions;
\item invariant subspaces or spectral projectors for separated clusters;
\item exact labels only with a separation gap.
\end{itemize}
If disks around roots are pairwise separated by \(\eta\), the least-disk
ordering is stable under perturbations smaller than \(\eta/3\).  The gap is a
guard and the label-selection proof is finite.  Without it, the framework
retains the relational cloud.

\subsection{Metastability versus a modulus}

A metastability statement has the form
\[
  \forall k\,\forall g\,\exists N\le B(k,g)\,
  \forall m,n\in[N,N+g(N)]\quad |q_m-q_n|\le2^{-k}.
\]
For each fixed \(k,g\), this is a finite-search theorem and fits the framework.
It does not supply a tail modulus valid for every \(m,n\ge N\).  Therefore it
supports finite-window conclusions but not, by itself, a regular name or
operations such as guarded reciprocal.  The type of the interface must retain
this distinction.
